% -----------------------------------------------------------------------------------------------------
\section{振幅的共线因子化与 Coulomb 相位}\label{sec:sudakov-collinear}

本节讨论振幅层次上两个与运动学组态密切相关的现象。其一是 Coulomb 相位：同为入射或同为出射的两条有色线之间的软交换带有 $\ii\pi$，它在类空与类时组态之间不同；
其二是共线因子化：当所有共线方向都在末态时，振幅在全阶上分解为硬、软、共线因子之积，而当一个入射部分子与一个出射部分子共线时，严格的共线因子化在单圈即失效。

\begin{proposition}[Coulomb 相位 \tagR]\label{prop:sudakov-coulomb}
（$R_{\mathrm{PT}}$：在重整化微扰论中对红外正规化的部分子态严格。）在单圈阶：
(a) 类空形状因子是实的，类时形状因子的虚部为
\begin{equation}\label{eq:sudakov-coulomb-ff}
\Imag F^{(1)}(s+\ii0)=\frac{\alpha_sC_F}{4\pi}\Big(\frac{\mu^2}{s}\Big)^{\varepsilon}\sin(\pi\varepsilon)\,f(\varepsilon)
=\frac{\alpha_sC_F}{4\pi}\Big[-\frac{2\pi}{\varepsilon}+2\pi\ell-3\pi\Big]+O(\varepsilon),\qquad \ell=\ln\frac{s}{\mu^2},
% VF-EE-4
\end{equation}
它在 $|F|^2$ 的 $O(\alpha_s)$ 项中不出现。
(b) 两条类时 Wilson 线在原点相遇（$\cosh\beta=v_1\cdot v_2$，$v_i^0>0$）时，单圈尖点反常维度在散射组态（一入一出，$W[\infty v_1,0]W[0,-\infty v_2]$）为实数
$\frac{\alpha_sC_F}\pi(\beta\coth\beta-1)$；在产生组态（同为出射，$W[\infty v_1,0]W[0,\infty v_2]$）与吸收组态（同为入射）为
\begin{equation}\label{eq:sudakov-coulomb-cusp}
\Gamma_{\rm cusp}^{\rm prod}(\beta)=\frac{\alpha_sC_F}{\pi}\big[(\beta-\ii\pi)\coth\beta-1\big],
\end{equation}
即 $\beta\to\beta-\ii\pi$；其虚部 $-\ii\alpha_sC_F\coth\beta$ 在类光极限下趋于 $-\ii\alpha_sC_F$。
(c) 因此在 $n$ 条 Wilson 线的单圈关联函数中（两两单胶子交换之和，第 $(i,j)$ 对的色因子为 $\mathbf T_i\cdot\mathbf T_j$），$\ii\pi$ 恰好出现在 $\sigma_{ij}=+1$ 的那些对上，
与式~\eqref{eq:sudakov-dipole} 中 $\ln\frac{\mu^2}{-s_{ij}}$ 的虚部 $\ii\pi$（$s_{ij}>0$ 时）一致。（本命题只涉及部分子态与 Wilson 线算符。）
\uses{prop:sudakov-oneloop,prop:appB-cusp-oneloop}
\end{proposition}

\begin{proof}
(a) 由\cref{prop:sudakov-oneloop}(b)，$F^{(1)}(s+\ii0)=\frac{\alpha_sC_F}{4\pi}(\mu^2/s)^\varepsilon\ee^{\ii\pi\varepsilon}f(\varepsilon)$，取虚部即得第一式；
$\sin(\pi\varepsilon)=\pi\varepsilon+O(\varepsilon^3)$、$(\mu^2/s)^\varepsilon=1-\varepsilon\ell+O(\varepsilon^2)$ 与 $f=-\frac2{\varepsilon^2}-\frac3\varepsilon+O(1)$ 相乘，
$\pi\varepsilon\cdot(-\frac2{\varepsilon^2})=-\frac{2\pi}\varepsilon$，$\pi\varepsilon\cdot[-\frac3\varepsilon+(-\varepsilon\ell)(-\frac2{\varepsilon^2})]=-3\pi+2\pi\ell$，即第二式。
类空区 $(\mu^2/Q^2)^\varepsilon f(\varepsilon)$ 为实数。$|F|^2=1+2\Real F^{(1)}+O(\alpha_s^2)$，虚部不进入 $O(\alpha_s)$。

(b) 散射组态即\cref{prop:appB-cusp-oneloop}。产生组态与它只有两处不同：$W[0,\infty v_2]$ 的路径取向为 $-v_2$，顶点多一个负号；两个顶点的间隔由 $lv_1+sv_2$ 变为 $lv_1-sv_2$（$l,s>0$）。
沿用附录~B 的记号（$l=\rho u$，$s=\rho(1-u)$，绝热阻尼 $\ee^{-\delta(l+s)}$），交换图的紫外极点变为
$-\frac{\alpha_sC_F}{2\pi\varepsilon}\cosh\beta\int_0^1\frac{\dd u}{-N'(u)+\ii0}$，其中 $(lv_1-sv_2)^2=\rho^2N'(u)$，
$N'(u)=1-2u(1-u)(1+\cosh\beta)$；$+\ii0$ 来自 Feynman 传播子 $(-x^2+\ii0)^{-1+\varepsilon}$。与散射组态不同，$N'$ 在 $(0,1)$ 内变号：两个顶点既可以类空分离也可以类时分离。
令 $w=2u-1$，$N'=\frac12(a^2w^2-b^2)$，$a^2=1+\cosh\beta=2\cosh^2\frac\beta2$，$b^2=\cosh\beta-1=2\sinh^2\frac\beta2$，$ab=\sinh\beta$；用 $\frac1{x-\ii0}=\mathrm{PV}\frac1x+\ii\pi\delta(x)$，
\[
\int_0^1\frac{\dd u}{N'(u)-\ii0}=\int_{-1}^{1}\frac{\dd w}{a^2w^2-b^2-\ii0}
=\frac{1}{ab}\ln\frac{a-b}{a+b}+\frac{\ii\pi}{ab}=\frac{\ii\pi-\beta}{\sinh\beta},
\]
其中用了 $\frac{a+b}{a-b}=\frac{\cosh(\beta/2)+\sinh(\beta/2)}{\cosh(\beta/2)-\sinh(\beta/2)}=\ee^{\beta}$，$\delta(a^2w^2-b^2)$ 在 $w=\pm b/a$ 处各贡献 $1/(2ab)$。
于是交换图的极点为 $-\frac{\alpha_sC_F}{2\pi\varepsilon}(\beta-\ii\pi)\coth\beta$；两条线的自能与组态无关，仍为 $+\frac{\alpha_sC_F}{2\pi\varepsilon}$（附录~B）。合起来
$W^{\rm prod}=1-\frac{\alpha_sC_F}{2\pi\varepsilon}[(\beta-\ii\pi)\coth\beta-1]+(\text{有限})$，按\cref{def:sudakov-cusp}的定义即式~\eqref{eq:sudakov-coulomb-cusp}。
吸收组态（同为入射，两个顶点位于 $-lv_1$、$-sv_2$）的间隔为 $-(lv_1-sv_2)$，平方不变，两个顶点的符号之积也与产生组态相同，结果相同。
由 $\coth(\beta-\ii\pi)=\coth\beta$，这正是在散射组态的结果中作代换 $\beta\to\beta-\ii\pi$（$\cosh(\beta-\ii\pi)=-\cosh\beta$ 对应于把一条线的方向反转）。

(c) 单圈时 $n$ 条线的关联函数是两两交换图与各线自能之和；第 $(i,j)$ 对交换图的色因子在色空间形式中为 $\mathbf T_i\cdot\mathbf T_j$（对色单态的夸克—反夸克对，$\mathbf T_1\cdot\mathbf T_2=-C_F$，
即上面的 $C_F$），其运动学因子按 (b) 只在两线同为入射或同为出射时含 $\ii\pi$。自能是实的且色对角。$\sigma_{ij}=+1$ 时 $s_{ij}=2p_i\cdot p_j>0$，
$\ln\frac{\mu^2}{-s_{ij}-\ii0}=\ln\frac{\mu^2}{s_{ij}}+\ii\pi$，与此相合。
\end{proof}

\cref{prop:sudakov-coulomb}说明了一个常被模糊表述的事实：形状因子的 $\ii\pi$ 与 $\pi^2$ 来自同为出射的两条线之间的软（Coulomb/Glauber 型）交换；在区域分析中，
加上解析正规子后 Glauber 区的积分为零（\cref{tab:sudakov-regions}），这些相位由带正确 $\ii0$ 方向的程函（Wilson 线）区再现。软函数的紫外极点虚部 $+\ii\frac{\alpha_sC_F}{2\varepsilon}$
（由式~\eqref{eq:sudakov-coulomb-cusp} 的类光极限）与形状因子红外极点的虚部 $-\ii\frac{\alpha_sC_F}{2\varepsilon}$（式~\eqref{eq:sudakov-coulomb-ff}）大小相等、符号相反，
正是\cref{ex:sudakov-onshell-hard}所述的“红外极点 = 无标度低能矩阵元的紫外极点”的又一例证（脚本 \texttt{s5\_cusp.py} 与 \texttt{s3\_irpoles.py}）。

\begin{proposition}[振幅的硬—软—共线因子化（出射或阈值运动学） \tagA]\label{prop:sudakov-amplitude-factorization}
设一个 QCD 散射振幅的外态只含沿若干方向的共线粒子与软粒子，且没有入射部分子与出射部分子共线（例如全部共线方向都在末态的衰变型过程，或部分子质心能量接近总能量、旁观者必为软的阈值运动学）。
则在领头幂次、对所有阶，振幅等于由 QCD 计算的若干算符矩阵元之积：每个共线方向一个共线矩阵元（夸克或胶子场乘以 Wilson 线），一个软矩阵元（沿各共线方向的 Wilson 线），
软—共线重叠由沿两个方向的折线 Wilson 线的真空矩阵元除去，系数只依赖于硬运动学；这一等式与正规子、规范无关，并蕴含软流与共线分裂函数对所有阶与过程无关
\cite{FeigeSchwartz2014,Feige:2013zla}。入射部分子与出射部分子共线的构型被明确排除在外\cite[第~2.2、10、11.2~节]{FeigeSchwartz2014}。（本命题只涉及部分子振幅。）
\uses{thm:sudakov-factorization,thm:ir-leading-regions}
\end{proposition}

\begin{proofsketch}
Feige 与 Schwartz 的证明分四步\cite[第~5--9~节]{FeigeSchwartz2014}：(1) “着色”，把图按软敏感性分解；(2) 用钳夹分析与幂次计数确定红外敏感区域（约化图），
其中只有末态共线方向时 Glauber 标度不产生领头奇点（由衰变率的红外有限性，同书第~2.2~节）；(3) 选取一种“因子化规范”，使软—共线解耦显明，鬼场退耦；
(4) 用 Ward 恒等式逐步自举出 Wilson 线，得到软—共线因子化。这些要素与\cref{thm:sudakov-factorization}相同（领头区域\cref{thm:ir-leading-regions}、Ward 恒等式、重叠的扣除），
区别在于外线数目任意、并在振幅层次上处理一切软与共线辐射。完整论证见\cite{FeigeSchwartz2014}。
\end{proofsketch}

\paragraph{地位与缺口}全阶论证（[A]）：这是一个已发表的、在振幅层次上的长篇全阶论证，本书没有独立重推；其最弱的步骤与\cref{thm:sudakov-factorization}相同——
领头区域与幂次计数只给出 E-S 型陈述，非阿贝尔 Ward 恒等式的组合学在其因子化规范中给出。适用范围受明确限制：入射—出射共线的构型被排除，因此一般的强子—强子散射（有高能旁观者）不在其内，
这与下面的类空失效相容。宽角振幅的红外森林公式（\cite{Ma2020}；第~24~章\cref{prop:front-ir-forest}，前向指针）给出了另一种全阶表述。

\begin{remark}[类空共线极限预览 \tagV]\label{rem:sudakov-spacelike}
当一个入射部分子与一个出射部分子变得共线（类空共线极限）时，单圈分裂振幅依赖于非共线部分子的色荷与动量：与过程无关的\emph{严格}共线因子化在单圈即失效，
而类时（全为出射）的共线极限中它成立\cite{CataniDeFlorianRodrigo2012}。这与\cref{prop:sudakov-amplitude-factorization}排除入射—出射共线构型相容。完整的陈述与讨论见第~18~章\cref{prop:break-spacelike}（前向指针）。
\uses{prop:sudakov-amplitude-factorization}
\end{remark}

\paragraph{地位与缺口}[V]，证据类型：显式的单圈（与两圈）计算\cite{CataniDeFlorianRodrigo2012}。由\cref{prop:sudakov-coulomb}(c) 可以看出单圈失效的来源：
类空共线对中的两个部分子一入一出，对每个非共线部分子 $j$，偶极子 $(1,j)$ 与 $(2,j)$ 中恰有一个 $\sigma=+1$，即恰有一个带 $\ii\pi$，哪一个取决于 $j$ 是入射还是出射；
于是 $\ii\pi$ 项不能合并为 $(\mathbf T_1+\mathbf T_2)\cdot\mathbf T_j$，奇异因子依赖于非共线部分子的色荷。它本身\emph{并不}说明 PDF 因子化失效：这是关于振幅的多部分子共线极限的陈述，它对单举截面的后果在许多情形中相消；
据 Catani、de Florian 与 Rodrigo 的摘要，它影响强子—强子截面中对数增强项的非阿贝尔结构（从 NNLO 起），并\emph{可能}影响质量奇异性的因子化（从 N$^3$LO 起）。
Drell--Yan 的共线因子化仍是全阶论证（[A]，\cref{ch:dy}）。

\paragraph{与前沿章节的联系。}本章的 cusp 结构在前沿章节中以另一种面貌出现：重—轻 Sudakov 形状因子与类空喷注函数（\cref{def:yzliu-hl-jet}），
以及离光锥软函数作为快速色源对的 HQET 型形状因子（\cref{def:lamet-soft-offlc}、\cref{thm:yzliu-soft-hqet}），其大角度常数正是\cref{sec:sudakov-evolution}中 $\beta\coth\beta-1$ 的 $-1$ 的全阶推广 $\gamma_{HH}$（均为前向指针，见\cref{ch:yzliu}）。

\begin{tcolorbox}[enhanced,breakable,title=常见误解,fonttitle=\bfseries,
  colback=gray!3!white,colframe=gray!55!black,boxrule=0.6pt,arc=2pt,
  left=5pt,right=5pt,top=3pt,bottom=3pt,before skip=6pt,after skip=6pt]
\textbf{「快度发散是坏正规子的人为产物。」}不对。有质量交换的形状因子中，单个共线区域的快度发散在纯维数正规化下无法正规化，大对数 $\ln(Q^2/M^2)$ 正是通过解析正规子极点的相消重新出现的
（\cref{ex:sudakov-massive}）；快度发散标志着一个新的因子化变量（快度或 $\zeta$），其依赖由演化方程控制（\cref{prop:sudakov-kg}；第~14~章\cref{prop:tmd-evolution}）。
不同正规子给出的单个区域不同，正确组合的物理量相同。\par
\textbf{「有质量的 Sudakov 形状因子必须用解析正规子。」}不准确。只有当\emph{交换的玻色子}有质量时才需要；在壳有质量夸克与无质量胶子的情形，共线区的软端 $\int\dd x\,x^{-1-2\varepsilon}$ 已被维数正规化调节，
硬区与两个共线区之和再现精确结果的领头幂次（\cref{ex:sudakov-massive}）。\par
\textbf{「类空严格共线因子化的失效说明 PDF 因子化失效。」}不对。那是振幅的类空多部分子共线极限的性质，对单举截面的后果在许多情形中相消；Drell--Yan 的共线因子化仍为 [A]
（\cref{rem:sudakov-spacelike}；\cref{ch:break}）。\par
\textbf{「偶极子公式给出红外极点的全阶精确形式。」}只在两圈以内成立；三圈起出现非偶极子项（\cref{prop:sudakov-dipole-violation}，[V]）。
\end{tcolorbox}

\paragraph{本章小结}夸克形状因子是最简单的全阶因子化定理的载体。一圈结果（\cref{prop:sudakov-oneloop}）对 $\varepsilon$ 精确，其双极点来自既软又共线的胶子，
向类时区的延拓带来约 $25\%$ 的 $\pi^2$ 修正；区域分析显示了硬、共线、超软三类标度及其极点的相消（\cref{ex:sudakov-onshell-hard,ex:sudakov-offshell}），并揭示了有质量交换时快度发散的必然性（\cref{ex:sudakov-massive}）。
全阶上，领头区域、共线与软近似、Ward 恒等式与减除给出 $F=H\,J_1J_2/S$（\cref{thm:sudakov-factorization}，[A]），其中 $J_i$ 与 $S$ 是含 Wilson 线的算符矩阵元；
对辅助快度的依赖导出 $K+G$ 方程（\cref{prop:sudakov-kg}），其反常维度是 cusp 反常维度的两倍，积分后红外极点与双对数指数化（\cref{cor:sudakov-exponentiation}）。
推广到多腿振幅，红外极点由色空间中的 $\mathbf Z$ 因子生成，两圈以内是偶极子形式，三圈起不是（\cref{thm:sudakov-ir-poles,prop:sudakov-dipole-violation}）。
本章 $[A]$ 结论的共同瓶颈是减除的积分范数估计（\cref{asm:appN-collins}，指针）与非阿贝尔多胶子 Ward 恒等式。

\paragraph{这些工具在后文的用途}一圈形状因子是 $e^+e^-$ 总截面的虚修正（第~6~章\cref{prop:ee-oneloop}）、Drell--Yan 系数函数中 $\delta(1-z)$ 与 $\pi^2$ 项的来源（第~12~章\cref{prop:dy-nlo}），
也是 SCET 中硬匹配系数 $C_V$ 的单圈值（第~15~章\cref{prop:scet-hard-matching}）；区域分析与离壳三角形在第~9~章\cref{def:method-regions}、\cref{ex:method-sudakov-subtraction}
与附录~N 的\cref{prop:appN-triangle-exact}中成为减除法的标准例子；Ward 恒等式把纵向胶子收集为 Wilson 线的机制在第~8~章的 PDF 定义（\cref{def:pdf-quark}）与第~9~章\cref{lem:method-ward-multi}中推广；
$K+G$ 方程与 cusp 反常维度是 Collins--Soper 演化（第~14~章\cref{prop:tmd-evolution}）与阈值重求和（第~13~章\cref{thm:thr-refactorization}）的原型；
多腿红外极点与 Coulomb 相位在喷注与超领头对数（第~16~章\cref{prop:jets-sll}）以及类空共线失效（第~18~章\cref{prop:break-spacelike}）中再次出现；
cusp 及其大角度常数在\cref{ch:yzliu}的重—轻 Sudakov 问题中出现。
